% =====================================================================
% Spacecraft Attitude Dynamics and Control Model — physics summary
%
% USAGE
%   In your Overleaf project's main .tex file, add:
%       % =====================================================================
% Spacecraft Attitude Dynamics and Control Model — physics summary
%
% USAGE
%   In your Overleaf project's main .tex file, add:
%       % =====================================================================
% Spacecraft Attitude Dynamics and Control Model — physics summary
%
% USAGE
%   In your Overleaf project's main .tex file, add:
%       % =====================================================================
% Spacecraft Attitude Dynamics and Control Model — physics summary
%
% USAGE
%   In your Overleaf project's main .tex file, add:
%       \input{attitude_dynamics_model.tex}
%   wherever you want the section to appear.
%
% REQUIRED PACKAGES (most Overleaf templates already load these)
%       \usepackage{amsmath, amssymb}
%       \usepackage{bm}     % bold math symbols
%
% BIBLIOGRAPHY
%   Add the following entry to your .bib file so that \cite{markley2014}
%   resolves correctly:
%
%   @book{markley2014,
%     author    = {Markley, F. Landis and Crassidis, John L.},
%     title     = {Fundamentals of Spacecraft Attitude Determination and Control},
%     publisher = {Springer},
%     year      = {2014},
%     series    = {Space Technology Library},
%     volume    = {33},
%     address   = {New York, NY},
%     doi       = {10.1007/978-1-4939-0802-8}
%   }
%
%   If you are not using BibTeX, replace the \cite{markley2014} occurrences
%   below with a manual reference such as "Markley \& Crassidis (2014)".
% =====================================================================


\section{Spacecraft Attitude Dynamics and Control Model}
\label{sec:attitude-model}

This section summarises the equations of motion used in the closed-loop
attitude simulation. Sign conventions and equation numbering follow
Markley and Crassidis~\cite{markley2014} throughout.


% ---------------------------------------------------------------------
\subsection{Quaternion Convention}
\label{sec:quat-convention}

We adopt the \emph{scalar-last} unit quaternion
\begin{equation}
\bm{q} \;=\; \begin{bmatrix} \bm{q}_v \\ q_4 \end{bmatrix} \in \mathbb{R}^{4},
\qquad \bm{q}_v \in \mathbb{R}^{3},\; q_4 \in \mathbb{R},\; \|\bm{q}\| = 1,
\label{eq:scalar-last-q}
\end{equation}
matching the layout used in Eq.~(2.84) of~\cite{markley2014}. The
quaternion product is defined such that
$\bm{A}(\bm{q}\otimes\bm{p}) = \bm{A}(\bm{q})\,\bm{A}(\bm{p})$
(Eq.~2.82b of~\cite{markley2014}):
\begin{equation}
\bm{q}\otimes\bm{p} \;=\;
\begin{bmatrix}
 q_4\,\bm{p}_v + p_4\,\bm{q}_v - \bm{q}_v\times\bm{p}_v \\[2pt]
 q_4\,p_4 - \bm{q}_v^{\!\top}\bm{p}_v
\end{bmatrix},
\qquad
\bm{q}^{-1} \;=\; \begin{bmatrix} -\bm{q}_v \\ q_4 \end{bmatrix}.
\end{equation}


% ---------------------------------------------------------------------
\subsection{Quaternion Kinematics}
\label{sec:quat-kinematics}

The rate of change of the body-to-inertial quaternion under body angular
velocity $\bm{\omega}\in\mathbb{R}^{3}$ is
(Eq.~2.88 of~\cite{markley2014})
\begin{equation}
\dot{\bm{q}}
\;=\; \tfrac{1}{2}\,\bm{\Omega}(\bm{\omega})\,\bm{q}
\;=\;
\begin{bmatrix}
 \tfrac{1}{2}\bigl(q_4\,\bm{\omega} - \bm{\omega}\times\bm{q}_v\bigr) \\[3pt]
-\tfrac{1}{2}\,\bm{\omega}^{\!\top}\bm{q}_v
\end{bmatrix}.
\label{eq:quat-kin}
\end{equation}


% ---------------------------------------------------------------------
\subsection{Rigid-Body Dynamics with Reaction Wheels}
\label{sec:euler-wheels}

For a spacecraft with symmetric, positive-definite inertia tensor
$\bm{J}\in\mathbb{R}^{3\times3}$ (Eq.~3.13 of~\cite{markley2014})
and internal wheel angular momentum $\bm{h}_w^{\,b}$ expressed in the body
frame, Euler's rotational equation generalises to (Sec.~7.2
of~\cite{markley2014})
\begin{equation}
\boxed{\;
\bm{J}\,\dot{\bm{\omega}}
\;=\;
\bm{u}_{\text{body}}
\;-\;
\bm{\omega}\times\!\bigl(\bm{J}\bm{\omega} + \bm{h}_w^{\,b}\bigr).
\;}
\label{eq:euler-wheels}
\end{equation}
The total external torque on the body collects three contributions:
\begin{equation}
\bm{u}_{\text{body}} \;=\; -\bm{W}\bm{\tau}_w \;+\; \bm{\tau}_{\text{mag}}
\;+\; \bm{T}_{\!s},
\label{eq:u-body-decomp}
\end{equation}
the wheel reaction torque, the magnetic-torquer torque, and the propellant
slosh torque defined in Sections~\ref{sec:bdot} and~\ref{sec:slosh}.


% ---------------------------------------------------------------------
\subsection{Reaction Wheels and Torque Allocation}
\label{sec:wheels}

We model $n_w$ momentum wheels whose spin-axis unit vectors form the
columns of a configuration matrix
$\bm{W}\in\mathbb{R}^{3\times n_w}$. The scalar wheel momenta about
their spin axes are stacked into $\bm{h}_w\in\mathbb{R}^{n_w}$, so the
body-frame wheel momentum is
\begin{equation}
\bm{h}_w^{\,b} \;=\; \bm{W}\,\bm{h}_w.
\end{equation}
With wheel torque commands $\bm{\tau}_w\in\mathbb{R}^{n_w}$, the wheels
react $-\bm{W}\bm{\tau}_w$ onto the body and their internal momentum
evolves as
\begin{equation}
\dot{\bm{h}}_w \;=\; \bm{\tau}_w.
\end{equation}
A desired body control torque $\bm{u}_{\text{cmd}}$ from the attitude
controller is mapped to wheel commands by the minimum-norm pseudoinverse
allocation
\begin{equation}
\bm{\tau}_w \;=\; -\bm{W}^{+}\,\bm{u}_{\text{cmd}},
\qquad
\bm{W}^{+} \;=\; \bm{W}^{\!\top}\!\bigl(\bm{W}\bm{W}^{\!\top}\bigr)^{-1}.
\label{eq:wheel-alloc}
\end{equation}

\paragraph{NASA standard four-wheel pyramid.}
The simulation uses the symmetric four-wheel pyramid configuration
(Sec.~7.2 of~\cite{markley2014}). Each spin axis is tilted by
$\beta = \arccos(1/\sqrt{3}) \approx 54.7356^\circ$ from the body $+z$
axis and the four wheels are spaced $90^\circ$ apart in azimuth. This
gives equal projection onto all three body axes and provides
single-fault tolerance.


% ---------------------------------------------------------------------
\subsection{Quaternion-Feedback PD Controller}
\label{sec:pd-controller}

Given a desired attitude quaternion $\bm{q}_{\text{des}}$, the attitude
error quaternion is
\begin{equation}
\delta\bm{q} \;=\; \bm{q}_{\text{des}}^{\,-1} \otimes \bm{q}
\;=\; \bigl[\,\delta\bm{q}_v;\; \delta q_4\,\bigr]^{\!\top}.
\end{equation}
The quaternion-feedback proportional--derivative law of Sec.~7.4
of~\cite{markley2014} is used, with element-wise saturation:
\begin{equation}
\bm{u}_{\text{cmd}} \;=\;
\operatorname{sat}_{u_{\max}}\!\Bigl(
 -K_p\,\operatorname{sign}(\delta q_4)\,\delta\bm{q}_v
 \;-\; K_d\,\bm{\omega}
\Bigr),
\label{eq:pd}
\end{equation}
where $K_p, K_d > 0$ are scalar gains and
$\operatorname{sat}_{u_{\max}}(\cdot)$ clips each component to
$[-u_{\max},\, +u_{\max}]$. The factor
$\operatorname{sign}(\delta q_4)$ selects the shortest rotation and
prevents the well-known \emph{unwinding} phenomenon.


% ---------------------------------------------------------------------
\subsection{Magnetic Momentum Unloading}
\label{sec:bdot}

Reaction wheels accumulate angular momentum from disturbance torques
until they saturate. Magnetic torquers transfer that momentum from the
wheels to the geomagnetic field (Sec.~7.5 of~\cite{markley2014}). A
magnetic dipole moment $\bm{m}\in\mathbb{R}^{3}$ commanded in a field
$\bm{B}\in\mathbb{R}^{3}$ produces a body torque
\begin{equation}
\bm{\tau}_{\text{mag}} \;=\; \bm{m}\times\bm{B}.
\end{equation}
We use the cross-product (``$\bm{B}$-dot-style'') unloading law
\begin{equation}
\bm{m} \;=\;
\operatorname{sat}_{m_{\max}}\!\biggl(
\frac{K_{\text{dump}}}{\|\bm{B}\|^{2}}\,
 \bm{h}_w^{\,b} \times \bm{B}
\biggr),
\label{eq:bdot}
\end{equation}
which (before saturation) yields the body torque
\begin{equation}
\bm{\tau}_{\text{mag}}
\;=\;
-K_{\text{dump}}\,
\Bigl(\bm{h}_w^{\,b} - (\bm{h}_w^{\,b}\!\cdot\!\hat{\bm{B}})\,\hat{\bm{B}}\Bigr),
\end{equation}
i.e.\ a damping torque applied to the component of $\bm{h}_w^{\,b}$
\emph{perpendicular} to $\bm{B}$. Orbital variation of $\bm{B}$ rotates
the unloadable subspace and sweeps it through all three body axes over
each orbit. The commanded dipole is saturated element-wise to
$\bm{m}\in[-m_{\max},\,+m_{\max}]^{3}$.

For the geomagnetic field, the simulation uses a simple inclined rotating
dipole of magnitude $\sim\!30\,\mu\mathrm{T}$ as a low-Earth-orbit proxy;
in flight-grade work this is replaced by an IGRF model rotated into the
body frame.


% ---------------------------------------------------------------------
\subsection{Propellant Slosh Torque}
\label{sec:slosh}

Liquid propellant slosh is modelled as a forced, damped second-order
oscillator acting on a 3-D torque state $\bm{T}_{\!s}\in\mathbb{R}^{3}$,
driven by body rate $\bm{\omega}$ and body angular acceleration
$\dot{\bm{\omega}}$:
\begin{equation}
\boxed{\;
\ddot{\bm{T}}_{\!s}
\;=\;
 -A_s(t)\,\bm{\omega}
 \;-\; B_s(t)\,\dot{\bm{\omega}}
 \;-\; C_s(t)\,\dot{\bm{T}}_{\!s}
 \;-\; K_s(t)\,\bm{T}_{\!s}.
\;}
\label{eq:slosh}
\end{equation}
The coefficients $A_s, B_s, C_s, K_s$ are scalars (identical on all three
axes) and admit an optional sinusoidal time variation about a positive
DC offset:
\begin{equation}
f(t) \;=\; f_{\text{nom}} + f_{\text{amp}}\sin(\omega_f\,t),
\qquad f \in \{A_s,\,B_s,\,C_s,\,K_s\}.
\label{eq:slosh-coeff}
\end{equation}
Requiring $f_{\text{nom}} > |f_{\text{amp}}|$ keeps $K_s,C_s$ strictly
positive (a negative stiffness or damping would be unphysical).

The slosh torque enters Euler's equation~\eqref{eq:euler-wheels} as an
external disturbance through $\bm{T}_{\!s}$ in~\eqref{eq:u-body-decomp}.
Because $\bm{T}_{\!s}$ is a \emph{state} (not an instantaneous algebraic
function of $\dot{\bm{\omega}}$), there is no implicit loop:
$\dot{\bm{\omega}}$ is computed from~\eqref{eq:euler-wheels} using the
current $\bm{T}_{\!s}$, and the just-computed $\dot{\bm{\omega}}$ is then
substituted into~\eqref{eq:slosh}.


% ---------------------------------------------------------------------
\subsection{State Vector and Equations of Motion}
\label{sec:state}

Collecting the previous results, the full simulation state is
\begin{equation}
\bm{x}(t) \;=\;
\bigl[\,\bm{q};\; \bm{\omega};\; \bm{h}_w;\; \bm{T}_{\!s};\;
       \dot{\bm{T}}_{\!s}\,\bigr]
\;\in\; \mathbb{R}^{\,13+n_w},
\end{equation}
governed by
\begin{align}
\dot{\bm{q}} &= \tfrac{1}{2}\,\bm{\Omega}(\bm{\omega})\,\bm{q}
   && \text{(kinematics, Eq.~\ref{eq:quat-kin})} \\[2pt]
\bm{J}\,\dot{\bm{\omega}} &= \bm{u}_{\text{body}}
 \;-\; \bm{\omega}\times\!\bigl(\bm{J}\bm{\omega}+\bm{W}\bm{h}_w\bigr)
   && \text{(Euler, Eq.~\ref{eq:euler-wheels})} \\[2pt]
\dot{\bm{h}}_w &= \bm{\tau}_w \;=\; -\bm{W}^{+}\bm{u}_{\text{cmd}}
   && \text{(wheel momentum)} \\[2pt]
\ddot{\bm{T}}_{\!s} &=
   -A_s\bm{\omega} - B_s\dot{\bm{\omega}}
   - C_s\dot{\bm{T}}_{\!s} - K_s\bm{T}_{\!s}
   && \text{(slosh, Eq.~\ref{eq:slosh})}
\end{align}
with
\(
\bm{u}_{\text{body}} \;=\;
-\bm{W}\bm{\tau}_w \;+\; \bm{m}\times\bm{B} \;+\; \bm{T}_{\!s},
\)
$\bm{u}_{\text{cmd}}$ given by~\eqref{eq:pd}, and $\bm{m}$
given by~\eqref{eq:bdot}.


% ---------------------------------------------------------------------
\subsection{Euler-Angle Conversion}
\label{sec:euler-conv}

For initial-condition specification and post-processing visualisation,
3-2-1 ($Z$-$Y$-$X$) yaw--pitch--roll Euler angles $(\phi,\theta,\psi)$
are converted to and from the scalar-last quaternion using
App.~B of~\cite{markley2014}:
\begin{equation}
\bm{q} \;=\; \bm{q}_z(\psi)\otimes\bm{q}_y(\theta)\otimes\bm{q}_x(\phi).
\end{equation}
The inverse extraction, used for the angle/rate output plots, is
\begin{align}
\phi   &= \operatorname{atan2}\!\bigl(2(q_4 q_1 + q_2 q_3),\;
                                      1 - 2(q_1^{2} + q_2^{2})\bigr), \\[2pt]
\theta &= \arcsin\!\bigl(\operatorname{clip}\bigl(2(q_4 q_2 - q_1 q_3),\,
                                                  -1,\,+1\bigr)\bigr), \\[2pt]
\psi   &= \operatorname{atan2}\!\bigl(2(q_4 q_3 + q_1 q_2),\;
                                      1 - 2(q_2^{2} + q_3^{2})\bigr),
\end{align}
where the $\arcsin$ argument is clipped to $[-1,+1]$ for numerical
robustness in the neighbourhood of the gimbal-lock singularity at
$\theta = \pm 90^\circ$.


% ---------------------------------------------------------------------
\subsection{Numerical Integration}
\label{sec:integration}

Two integrators are provided for the augmented state:
\begin{itemize}\setlength{\itemsep}{2pt}
  \item a hand-written fixed-step classical fourth-order Runge--Kutta
        (RK4) with quaternion re-normalisation after each step
        (zero-order hold on controller outputs over each $\Delta t$);
  \item an adaptive solver via SciPy's
        \texttt{scipy.integrate.solve\_ivp}
        (\texttt{RK45} / \texttt{DOP853} / \texttt{LSODA}), with the
        closed-loop treated as a continuous-time system and quaternion
        re-normalisation applied in post-processing.
\end{itemize}
Both integrators share the same dynamics, controller, and slosh
modules; only the time-stepping strategy differs.

%   wherever you want the section to appear.
%
% REQUIRED PACKAGES (most Overleaf templates already load these)
%       \usepackage{amsmath, amssymb}
%       \usepackage{bm}     % bold math symbols
%
% BIBLIOGRAPHY
%   Add the following entry to your .bib file so that \cite{markley2014}
%   resolves correctly:
%
%   @book{markley2014,
%     author    = {Markley, F. Landis and Crassidis, John L.},
%     title     = {Fundamentals of Spacecraft Attitude Determination and Control},
%     publisher = {Springer},
%     year      = {2014},
%     series    = {Space Technology Library},
%     volume    = {33},
%     address   = {New York, NY},
%     doi       = {10.1007/978-1-4939-0802-8}
%   }
%
%   If you are not using BibTeX, replace the \cite{markley2014} occurrences
%   below with a manual reference such as "Markley \& Crassidis (2014)".
% =====================================================================


\section{Spacecraft Attitude Dynamics and Control Model}
\label{sec:attitude-model}

This section summarises the equations of motion used in the closed-loop
attitude simulation. Sign conventions and equation numbering follow
Markley and Crassidis~\cite{markley2014} throughout.


% ---------------------------------------------------------------------
\subsection{Quaternion Convention}
\label{sec:quat-convention}

We adopt the \emph{scalar-last} unit quaternion
\begin{equation}
\bm{q} \;=\; \begin{bmatrix} \bm{q}_v \\ q_4 \end{bmatrix} \in \mathbb{R}^{4},
\qquad \bm{q}_v \in \mathbb{R}^{3},\; q_4 \in \mathbb{R},\; \|\bm{q}\| = 1,
\label{eq:scalar-last-q}
\end{equation}
matching the layout used in Eq.~(2.84) of~\cite{markley2014}. The
quaternion product is defined such that
$\bm{A}(\bm{q}\otimes\bm{p}) = \bm{A}(\bm{q})\,\bm{A}(\bm{p})$
(Eq.~2.82b of~\cite{markley2014}):
\begin{equation}
\bm{q}\otimes\bm{p} \;=\;
\begin{bmatrix}
 q_4\,\bm{p}_v + p_4\,\bm{q}_v - \bm{q}_v\times\bm{p}_v \\[2pt]
 q_4\,p_4 - \bm{q}_v^{\!\top}\bm{p}_v
\end{bmatrix},
\qquad
\bm{q}^{-1} \;=\; \begin{bmatrix} -\bm{q}_v \\ q_4 \end{bmatrix}.
\end{equation}


% ---------------------------------------------------------------------
\subsection{Quaternion Kinematics}
\label{sec:quat-kinematics}

The rate of change of the body-to-inertial quaternion under body angular
velocity $\bm{\omega}\in\mathbb{R}^{3}$ is
(Eq.~2.88 of~\cite{markley2014})
\begin{equation}
\dot{\bm{q}}
\;=\; \tfrac{1}{2}\,\bm{\Omega}(\bm{\omega})\,\bm{q}
\;=\;
\begin{bmatrix}
 \tfrac{1}{2}\bigl(q_4\,\bm{\omega} - \bm{\omega}\times\bm{q}_v\bigr) \\[3pt]
-\tfrac{1}{2}\,\bm{\omega}^{\!\top}\bm{q}_v
\end{bmatrix}.
\label{eq:quat-kin}
\end{equation}


% ---------------------------------------------------------------------
\subsection{Rigid-Body Dynamics with Reaction Wheels}
\label{sec:euler-wheels}

For a spacecraft with symmetric, positive-definite inertia tensor
$\bm{J}\in\mathbb{R}^{3\times3}$ (Eq.~3.13 of~\cite{markley2014})
and internal wheel angular momentum $\bm{h}_w^{\,b}$ expressed in the body
frame, Euler's rotational equation generalises to (Sec.~7.2
of~\cite{markley2014})
\begin{equation}
\boxed{\;
\bm{J}\,\dot{\bm{\omega}}
\;=\;
\bm{u}_{\text{body}}
\;-\;
\bm{\omega}\times\!\bigl(\bm{J}\bm{\omega} + \bm{h}_w^{\,b}\bigr).
\;}
\label{eq:euler-wheels}
\end{equation}
The total external torque on the body collects three contributions:
\begin{equation}
\bm{u}_{\text{body}} \;=\; -\bm{W}\bm{\tau}_w \;+\; \bm{\tau}_{\text{mag}}
\;+\; \bm{T}_{\!s},
\label{eq:u-body-decomp}
\end{equation}
the wheel reaction torque, the magnetic-torquer torque, and the propellant
slosh torque defined in Sections~\ref{sec:bdot} and~\ref{sec:slosh}.


% ---------------------------------------------------------------------
\subsection{Reaction Wheels and Torque Allocation}
\label{sec:wheels}

We model $n_w$ momentum wheels whose spin-axis unit vectors form the
columns of a configuration matrix
$\bm{W}\in\mathbb{R}^{3\times n_w}$. The scalar wheel momenta about
their spin axes are stacked into $\bm{h}_w\in\mathbb{R}^{n_w}$, so the
body-frame wheel momentum is
\begin{equation}
\bm{h}_w^{\,b} \;=\; \bm{W}\,\bm{h}_w.
\end{equation}
With wheel torque commands $\bm{\tau}_w\in\mathbb{R}^{n_w}$, the wheels
react $-\bm{W}\bm{\tau}_w$ onto the body and their internal momentum
evolves as
\begin{equation}
\dot{\bm{h}}_w \;=\; \bm{\tau}_w.
\end{equation}
A desired body control torque $\bm{u}_{\text{cmd}}$ from the attitude
controller is mapped to wheel commands by the minimum-norm pseudoinverse
allocation
\begin{equation}
\bm{\tau}_w \;=\; -\bm{W}^{+}\,\bm{u}_{\text{cmd}},
\qquad
\bm{W}^{+} \;=\; \bm{W}^{\!\top}\!\bigl(\bm{W}\bm{W}^{\!\top}\bigr)^{-1}.
\label{eq:wheel-alloc}
\end{equation}

\paragraph{NASA standard four-wheel pyramid.}
The simulation uses the symmetric four-wheel pyramid configuration
(Sec.~7.2 of~\cite{markley2014}). Each spin axis is tilted by
$\beta = \arccos(1/\sqrt{3}) \approx 54.7356^\circ$ from the body $+z$
axis and the four wheels are spaced $90^\circ$ apart in azimuth. This
gives equal projection onto all three body axes and provides
single-fault tolerance.


% ---------------------------------------------------------------------
\subsection{Quaternion-Feedback PD Controller}
\label{sec:pd-controller}

Given a desired attitude quaternion $\bm{q}_{\text{des}}$, the attitude
error quaternion is
\begin{equation}
\delta\bm{q} \;=\; \bm{q}_{\text{des}}^{\,-1} \otimes \bm{q}
\;=\; \bigl[\,\delta\bm{q}_v;\; \delta q_4\,\bigr]^{\!\top}.
\end{equation}
The quaternion-feedback proportional--derivative law of Sec.~7.4
of~\cite{markley2014} is used, with element-wise saturation:
\begin{equation}
\bm{u}_{\text{cmd}} \;=\;
\operatorname{sat}_{u_{\max}}\!\Bigl(
 -K_p\,\operatorname{sign}(\delta q_4)\,\delta\bm{q}_v
 \;-\; K_d\,\bm{\omega}
\Bigr),
\label{eq:pd}
\end{equation}
where $K_p, K_d > 0$ are scalar gains and
$\operatorname{sat}_{u_{\max}}(\cdot)$ clips each component to
$[-u_{\max},\, +u_{\max}]$. The factor
$\operatorname{sign}(\delta q_4)$ selects the shortest rotation and
prevents the well-known \emph{unwinding} phenomenon.


% ---------------------------------------------------------------------
\subsection{Magnetic Momentum Unloading}
\label{sec:bdot}

Reaction wheels accumulate angular momentum from disturbance torques
until they saturate. Magnetic torquers transfer that momentum from the
wheels to the geomagnetic field (Sec.~7.5 of~\cite{markley2014}). A
magnetic dipole moment $\bm{m}\in\mathbb{R}^{3}$ commanded in a field
$\bm{B}\in\mathbb{R}^{3}$ produces a body torque
\begin{equation}
\bm{\tau}_{\text{mag}} \;=\; \bm{m}\times\bm{B}.
\end{equation}
We use the cross-product (``$\bm{B}$-dot-style'') unloading law
\begin{equation}
\bm{m} \;=\;
\operatorname{sat}_{m_{\max}}\!\biggl(
\frac{K_{\text{dump}}}{\|\bm{B}\|^{2}}\,
 \bm{h}_w^{\,b} \times \bm{B}
\biggr),
\label{eq:bdot}
\end{equation}
which (before saturation) yields the body torque
\begin{equation}
\bm{\tau}_{\text{mag}}
\;=\;
-K_{\text{dump}}\,
\Bigl(\bm{h}_w^{\,b} - (\bm{h}_w^{\,b}\!\cdot\!\hat{\bm{B}})\,\hat{\bm{B}}\Bigr),
\end{equation}
i.e.\ a damping torque applied to the component of $\bm{h}_w^{\,b}$
\emph{perpendicular} to $\bm{B}$. Orbital variation of $\bm{B}$ rotates
the unloadable subspace and sweeps it through all three body axes over
each orbit. The commanded dipole is saturated element-wise to
$\bm{m}\in[-m_{\max},\,+m_{\max}]^{3}$.

For the geomagnetic field, the simulation uses a simple inclined rotating
dipole of magnitude $\sim\!30\,\mu\mathrm{T}$ as a low-Earth-orbit proxy;
in flight-grade work this is replaced by an IGRF model rotated into the
body frame.


% ---------------------------------------------------------------------
\subsection{Propellant Slosh Torque}
\label{sec:slosh}

Liquid propellant slosh is modelled as a forced, damped second-order
oscillator acting on a 3-D torque state $\bm{T}_{\!s}\in\mathbb{R}^{3}$,
driven by body rate $\bm{\omega}$ and body angular acceleration
$\dot{\bm{\omega}}$:
\begin{equation}
\boxed{\;
\ddot{\bm{T}}_{\!s}
\;=\;
 -A_s(t)\,\bm{\omega}
 \;-\; B_s(t)\,\dot{\bm{\omega}}
 \;-\; C_s(t)\,\dot{\bm{T}}_{\!s}
 \;-\; K_s(t)\,\bm{T}_{\!s}.
\;}
\label{eq:slosh}
\end{equation}
The coefficients $A_s, B_s, C_s, K_s$ are scalars (identical on all three
axes) and admit an optional sinusoidal time variation about a positive
DC offset:
\begin{equation}
f(t) \;=\; f_{\text{nom}} + f_{\text{amp}}\sin(\omega_f\,t),
\qquad f \in \{A_s,\,B_s,\,C_s,\,K_s\}.
\label{eq:slosh-coeff}
\end{equation}
Requiring $f_{\text{nom}} > |f_{\text{amp}}|$ keeps $K_s,C_s$ strictly
positive (a negative stiffness or damping would be unphysical).

The slosh torque enters Euler's equation~\eqref{eq:euler-wheels} as an
external disturbance through $\bm{T}_{\!s}$ in~\eqref{eq:u-body-decomp}.
Because $\bm{T}_{\!s}$ is a \emph{state} (not an instantaneous algebraic
function of $\dot{\bm{\omega}}$), there is no implicit loop:
$\dot{\bm{\omega}}$ is computed from~\eqref{eq:euler-wheels} using the
current $\bm{T}_{\!s}$, and the just-computed $\dot{\bm{\omega}}$ is then
substituted into~\eqref{eq:slosh}.


% ---------------------------------------------------------------------
\subsection{State Vector and Equations of Motion}
\label{sec:state}

Collecting the previous results, the full simulation state is
\begin{equation}
\bm{x}(t) \;=\;
\bigl[\,\bm{q};\; \bm{\omega};\; \bm{h}_w;\; \bm{T}_{\!s};\;
       \dot{\bm{T}}_{\!s}\,\bigr]
\;\in\; \mathbb{R}^{\,13+n_w},
\end{equation}
governed by
\begin{align}
\dot{\bm{q}} &= \tfrac{1}{2}\,\bm{\Omega}(\bm{\omega})\,\bm{q}
   && \text{(kinematics, Eq.~\ref{eq:quat-kin})} \\[2pt]
\bm{J}\,\dot{\bm{\omega}} &= \bm{u}_{\text{body}}
 \;-\; \bm{\omega}\times\!\bigl(\bm{J}\bm{\omega}+\bm{W}\bm{h}_w\bigr)
   && \text{(Euler, Eq.~\ref{eq:euler-wheels})} \\[2pt]
\dot{\bm{h}}_w &= \bm{\tau}_w \;=\; -\bm{W}^{+}\bm{u}_{\text{cmd}}
   && \text{(wheel momentum)} \\[2pt]
\ddot{\bm{T}}_{\!s} &=
   -A_s\bm{\omega} - B_s\dot{\bm{\omega}}
   - C_s\dot{\bm{T}}_{\!s} - K_s\bm{T}_{\!s}
   && \text{(slosh, Eq.~\ref{eq:slosh})}
\end{align}
with
\(
\bm{u}_{\text{body}} \;=\;
-\bm{W}\bm{\tau}_w \;+\; \bm{m}\times\bm{B} \;+\; \bm{T}_{\!s},
\)
$\bm{u}_{\text{cmd}}$ given by~\eqref{eq:pd}, and $\bm{m}$
given by~\eqref{eq:bdot}.


% ---------------------------------------------------------------------
\subsection{Euler-Angle Conversion}
\label{sec:euler-conv}

For initial-condition specification and post-processing visualisation,
3-2-1 ($Z$-$Y$-$X$) yaw--pitch--roll Euler angles $(\phi,\theta,\psi)$
are converted to and from the scalar-last quaternion using
App.~B of~\cite{markley2014}:
\begin{equation}
\bm{q} \;=\; \bm{q}_z(\psi)\otimes\bm{q}_y(\theta)\otimes\bm{q}_x(\phi).
\end{equation}
The inverse extraction, used for the angle/rate output plots, is
\begin{align}
\phi   &= \operatorname{atan2}\!\bigl(2(q_4 q_1 + q_2 q_3),\;
                                      1 - 2(q_1^{2} + q_2^{2})\bigr), \\[2pt]
\theta &= \arcsin\!\bigl(\operatorname{clip}\bigl(2(q_4 q_2 - q_1 q_3),\,
                                                  -1,\,+1\bigr)\bigr), \\[2pt]
\psi   &= \operatorname{atan2}\!\bigl(2(q_4 q_3 + q_1 q_2),\;
                                      1 - 2(q_2^{2} + q_3^{2})\bigr),
\end{align}
where the $\arcsin$ argument is clipped to $[-1,+1]$ for numerical
robustness in the neighbourhood of the gimbal-lock singularity at
$\theta = \pm 90^\circ$.


% ---------------------------------------------------------------------
\subsection{Numerical Integration}
\label{sec:integration}

Two integrators are provided for the augmented state:
\begin{itemize}\setlength{\itemsep}{2pt}
  \item a hand-written fixed-step classical fourth-order Runge--Kutta
        (RK4) with quaternion re-normalisation after each step
        (zero-order hold on controller outputs over each $\Delta t$);
  \item an adaptive solver via SciPy's
        \texttt{scipy.integrate.solve\_ivp}
        (\texttt{RK45} / \texttt{DOP853} / \texttt{LSODA}), with the
        closed-loop treated as a continuous-time system and quaternion
        re-normalisation applied in post-processing.
\end{itemize}
Both integrators share the same dynamics, controller, and slosh
modules; only the time-stepping strategy differs.

%   wherever you want the section to appear.
%
% REQUIRED PACKAGES (most Overleaf templates already load these)
%       \usepackage{amsmath, amssymb}
%       \usepackage{bm}     % bold math symbols
%
% BIBLIOGRAPHY
%   Add the following entry to your .bib file so that \cite{markley2014}
%   resolves correctly:
%
%   @book{markley2014,
%     author    = {Markley, F. Landis and Crassidis, John L.},
%     title     = {Fundamentals of Spacecraft Attitude Determination and Control},
%     publisher = {Springer},
%     year      = {2014},
%     series    = {Space Technology Library},
%     volume    = {33},
%     address   = {New York, NY},
%     doi       = {10.1007/978-1-4939-0802-8}
%   }
%
%   If you are not using BibTeX, replace the \cite{markley2014} occurrences
%   below with a manual reference such as "Markley \& Crassidis (2014)".
% =====================================================================


\section{Spacecraft Attitude Dynamics and Control Model}
\label{sec:attitude-model}

This section summarises the equations of motion used in the closed-loop
attitude simulation. Sign conventions and equation numbering follow
Markley and Crassidis~\cite{markley2014} throughout.


% ---------------------------------------------------------------------
\subsection{Quaternion Convention}
\label{sec:quat-convention}

We adopt the \emph{scalar-last} unit quaternion
\begin{equation}
\bm{q} \;=\; \begin{bmatrix} \bm{q}_v \\ q_4 \end{bmatrix} \in \mathbb{R}^{4},
\qquad \bm{q}_v \in \mathbb{R}^{3},\; q_4 \in \mathbb{R},\; \|\bm{q}\| = 1,
\label{eq:scalar-last-q}
\end{equation}
matching the layout used in Eq.~(2.84) of~\cite{markley2014}. The
quaternion product is defined such that
$\bm{A}(\bm{q}\otimes\bm{p}) = \bm{A}(\bm{q})\,\bm{A}(\bm{p})$
(Eq.~2.82b of~\cite{markley2014}):
\begin{equation}
\bm{q}\otimes\bm{p} \;=\;
\begin{bmatrix}
 q_4\,\bm{p}_v + p_4\,\bm{q}_v - \bm{q}_v\times\bm{p}_v \\[2pt]
 q_4\,p_4 - \bm{q}_v^{\!\top}\bm{p}_v
\end{bmatrix},
\qquad
\bm{q}^{-1} \;=\; \begin{bmatrix} -\bm{q}_v \\ q_4 \end{bmatrix}.
\end{equation}


% ---------------------------------------------------------------------
\subsection{Quaternion Kinematics}
\label{sec:quat-kinematics}

The rate of change of the body-to-inertial quaternion under body angular
velocity $\bm{\omega}\in\mathbb{R}^{3}$ is
(Eq.~2.88 of~\cite{markley2014})
\begin{equation}
\dot{\bm{q}}
\;=\; \tfrac{1}{2}\,\bm{\Omega}(\bm{\omega})\,\bm{q}
\;=\;
\begin{bmatrix}
 \tfrac{1}{2}\bigl(q_4\,\bm{\omega} - \bm{\omega}\times\bm{q}_v\bigr) \\[3pt]
-\tfrac{1}{2}\,\bm{\omega}^{\!\top}\bm{q}_v
\end{bmatrix}.
\label{eq:quat-kin}
\end{equation}


% ---------------------------------------------------------------------
\subsection{Rigid-Body Dynamics with Reaction Wheels}
\label{sec:euler-wheels}

For a spacecraft with symmetric, positive-definite inertia tensor
$\bm{J}\in\mathbb{R}^{3\times3}$ (Eq.~3.13 of~\cite{markley2014})
and internal wheel angular momentum $\bm{h}_w^{\,b}$ expressed in the body
frame, Euler's rotational equation generalises to (Sec.~7.2
of~\cite{markley2014})
\begin{equation}
\boxed{\;
\bm{J}\,\dot{\bm{\omega}}
\;=\;
\bm{u}_{\text{body}}
\;-\;
\bm{\omega}\times\!\bigl(\bm{J}\bm{\omega} + \bm{h}_w^{\,b}\bigr).
\;}
\label{eq:euler-wheels}
\end{equation}
The total external torque on the body collects three contributions:
\begin{equation}
\bm{u}_{\text{body}} \;=\; -\bm{W}\bm{\tau}_w \;+\; \bm{\tau}_{\text{mag}}
\;+\; \bm{T}_{\!s},
\label{eq:u-body-decomp}
\end{equation}
the wheel reaction torque, the magnetic-torquer torque, and the propellant
slosh torque defined in Sections~\ref{sec:bdot} and~\ref{sec:slosh}.


% ---------------------------------------------------------------------
\subsection{Reaction Wheels and Torque Allocation}
\label{sec:wheels}

We model $n_w$ momentum wheels whose spin-axis unit vectors form the
columns of a configuration matrix
$\bm{W}\in\mathbb{R}^{3\times n_w}$. The scalar wheel momenta about
their spin axes are stacked into $\bm{h}_w\in\mathbb{R}^{n_w}$, so the
body-frame wheel momentum is
\begin{equation}
\bm{h}_w^{\,b} \;=\; \bm{W}\,\bm{h}_w.
\end{equation}
With wheel torque commands $\bm{\tau}_w\in\mathbb{R}^{n_w}$, the wheels
react $-\bm{W}\bm{\tau}_w$ onto the body and their internal momentum
evolves as
\begin{equation}
\dot{\bm{h}}_w \;=\; \bm{\tau}_w.
\end{equation}
A desired body control torque $\bm{u}_{\text{cmd}}$ from the attitude
controller is mapped to wheel commands by the minimum-norm pseudoinverse
allocation
\begin{equation}
\bm{\tau}_w \;=\; -\bm{W}^{+}\,\bm{u}_{\text{cmd}},
\qquad
\bm{W}^{+} \;=\; \bm{W}^{\!\top}\!\bigl(\bm{W}\bm{W}^{\!\top}\bigr)^{-1}.
\label{eq:wheel-alloc}
\end{equation}

\paragraph{NASA standard four-wheel pyramid.}
The simulation uses the symmetric four-wheel pyramid configuration
(Sec.~7.2 of~\cite{markley2014}). Each spin axis is tilted by
$\beta = \arccos(1/\sqrt{3}) \approx 54.7356^\circ$ from the body $+z$
axis and the four wheels are spaced $90^\circ$ apart in azimuth. This
gives equal projection onto all three body axes and provides
single-fault tolerance.


% ---------------------------------------------------------------------
\subsection{Quaternion-Feedback PD Controller}
\label{sec:pd-controller}

Given a desired attitude quaternion $\bm{q}_{\text{des}}$, the attitude
error quaternion is
\begin{equation}
\delta\bm{q} \;=\; \bm{q}_{\text{des}}^{\,-1} \otimes \bm{q}
\;=\; \bigl[\,\delta\bm{q}_v;\; \delta q_4\,\bigr]^{\!\top}.
\end{equation}
The quaternion-feedback proportional--derivative law of Sec.~7.4
of~\cite{markley2014} is used, with element-wise saturation:
\begin{equation}
\bm{u}_{\text{cmd}} \;=\;
\operatorname{sat}_{u_{\max}}\!\Bigl(
 -K_p\,\operatorname{sign}(\delta q_4)\,\delta\bm{q}_v
 \;-\; K_d\,\bm{\omega}
\Bigr),
\label{eq:pd}
\end{equation}
where $K_p, K_d > 0$ are scalar gains and
$\operatorname{sat}_{u_{\max}}(\cdot)$ clips each component to
$[-u_{\max},\, +u_{\max}]$. The factor
$\operatorname{sign}(\delta q_4)$ selects the shortest rotation and
prevents the well-known \emph{unwinding} phenomenon.


% ---------------------------------------------------------------------
\subsection{Magnetic Momentum Unloading}
\label{sec:bdot}

Reaction wheels accumulate angular momentum from disturbance torques
until they saturate. Magnetic torquers transfer that momentum from the
wheels to the geomagnetic field (Sec.~7.5 of~\cite{markley2014}). A
magnetic dipole moment $\bm{m}\in\mathbb{R}^{3}$ commanded in a field
$\bm{B}\in\mathbb{R}^{3}$ produces a body torque
\begin{equation}
\bm{\tau}_{\text{mag}} \;=\; \bm{m}\times\bm{B}.
\end{equation}
We use the cross-product (``$\bm{B}$-dot-style'') unloading law
\begin{equation}
\bm{m} \;=\;
\operatorname{sat}_{m_{\max}}\!\biggl(
\frac{K_{\text{dump}}}{\|\bm{B}\|^{2}}\,
 \bm{h}_w^{\,b} \times \bm{B}
\biggr),
\label{eq:bdot}
\end{equation}
which (before saturation) yields the body torque
\begin{equation}
\bm{\tau}_{\text{mag}}
\;=\;
-K_{\text{dump}}\,
\Bigl(\bm{h}_w^{\,b} - (\bm{h}_w^{\,b}\!\cdot\!\hat{\bm{B}})\,\hat{\bm{B}}\Bigr),
\end{equation}
i.e.\ a damping torque applied to the component of $\bm{h}_w^{\,b}$
\emph{perpendicular} to $\bm{B}$. Orbital variation of $\bm{B}$ rotates
the unloadable subspace and sweeps it through all three body axes over
each orbit. The commanded dipole is saturated element-wise to
$\bm{m}\in[-m_{\max},\,+m_{\max}]^{3}$.

For the geomagnetic field, the simulation uses a simple inclined rotating
dipole of magnitude $\sim\!30\,\mu\mathrm{T}$ as a low-Earth-orbit proxy;
in flight-grade work this is replaced by an IGRF model rotated into the
body frame.


% ---------------------------------------------------------------------
\subsection{Propellant Slosh Torque}
\label{sec:slosh}

Liquid propellant slosh is modelled as a forced, damped second-order
oscillator acting on a 3-D torque state $\bm{T}_{\!s}\in\mathbb{R}^{3}$,
driven by body rate $\bm{\omega}$ and body angular acceleration
$\dot{\bm{\omega}}$:
\begin{equation}
\boxed{\;
\ddot{\bm{T}}_{\!s}
\;=\;
 -A_s(t)\,\bm{\omega}
 \;-\; B_s(t)\,\dot{\bm{\omega}}
 \;-\; C_s(t)\,\dot{\bm{T}}_{\!s}
 \;-\; K_s(t)\,\bm{T}_{\!s}.
\;}
\label{eq:slosh}
\end{equation}
The coefficients $A_s, B_s, C_s, K_s$ are scalars (identical on all three
axes) and admit an optional sinusoidal time variation about a positive
DC offset:
\begin{equation}
f(t) \;=\; f_{\text{nom}} + f_{\text{amp}}\sin(\omega_f\,t),
\qquad f \in \{A_s,\,B_s,\,C_s,\,K_s\}.
\label{eq:slosh-coeff}
\end{equation}
Requiring $f_{\text{nom}} > |f_{\text{amp}}|$ keeps $K_s,C_s$ strictly
positive (a negative stiffness or damping would be unphysical).

The slosh torque enters Euler's equation~\eqref{eq:euler-wheels} as an
external disturbance through $\bm{T}_{\!s}$ in~\eqref{eq:u-body-decomp}.
Because $\bm{T}_{\!s}$ is a \emph{state} (not an instantaneous algebraic
function of $\dot{\bm{\omega}}$), there is no implicit loop:
$\dot{\bm{\omega}}$ is computed from~\eqref{eq:euler-wheels} using the
current $\bm{T}_{\!s}$, and the just-computed $\dot{\bm{\omega}}$ is then
substituted into~\eqref{eq:slosh}.


% ---------------------------------------------------------------------
\subsection{State Vector and Equations of Motion}
\label{sec:state}

Collecting the previous results, the full simulation state is
\begin{equation}
\bm{x}(t) \;=\;
\bigl[\,\bm{q};\; \bm{\omega};\; \bm{h}_w;\; \bm{T}_{\!s};\;
       \dot{\bm{T}}_{\!s}\,\bigr]
\;\in\; \mathbb{R}^{\,13+n_w},
\end{equation}
governed by
\begin{align}
\dot{\bm{q}} &= \tfrac{1}{2}\,\bm{\Omega}(\bm{\omega})\,\bm{q}
   && \text{(kinematics, Eq.~\ref{eq:quat-kin})} \\[2pt]
\bm{J}\,\dot{\bm{\omega}} &= \bm{u}_{\text{body}}
 \;-\; \bm{\omega}\times\!\bigl(\bm{J}\bm{\omega}+\bm{W}\bm{h}_w\bigr)
   && \text{(Euler, Eq.~\ref{eq:euler-wheels})} \\[2pt]
\dot{\bm{h}}_w &= \bm{\tau}_w \;=\; -\bm{W}^{+}\bm{u}_{\text{cmd}}
   && \text{(wheel momentum)} \\[2pt]
\ddot{\bm{T}}_{\!s} &=
   -A_s\bm{\omega} - B_s\dot{\bm{\omega}}
   - C_s\dot{\bm{T}}_{\!s} - K_s\bm{T}_{\!s}
   && \text{(slosh, Eq.~\ref{eq:slosh})}
\end{align}
with
\(
\bm{u}_{\text{body}} \;=\;
-\bm{W}\bm{\tau}_w \;+\; \bm{m}\times\bm{B} \;+\; \bm{T}_{\!s},
\)
$\bm{u}_{\text{cmd}}$ given by~\eqref{eq:pd}, and $\bm{m}$
given by~\eqref{eq:bdot}.


% ---------------------------------------------------------------------
\subsection{Euler-Angle Conversion}
\label{sec:euler-conv}

For initial-condition specification and post-processing visualisation,
3-2-1 ($Z$-$Y$-$X$) yaw--pitch--roll Euler angles $(\phi,\theta,\psi)$
are converted to and from the scalar-last quaternion using
App.~B of~\cite{markley2014}:
\begin{equation}
\bm{q} \;=\; \bm{q}_z(\psi)\otimes\bm{q}_y(\theta)\otimes\bm{q}_x(\phi).
\end{equation}
The inverse extraction, used for the angle/rate output plots, is
\begin{align}
\phi   &= \operatorname{atan2}\!\bigl(2(q_4 q_1 + q_2 q_3),\;
                                      1 - 2(q_1^{2} + q_2^{2})\bigr), \\[2pt]
\theta &= \arcsin\!\bigl(\operatorname{clip}\bigl(2(q_4 q_2 - q_1 q_3),\,
                                                  -1,\,+1\bigr)\bigr), \\[2pt]
\psi   &= \operatorname{atan2}\!\bigl(2(q_4 q_3 + q_1 q_2),\;
                                      1 - 2(q_2^{2} + q_3^{2})\bigr),
\end{align}
where the $\arcsin$ argument is clipped to $[-1,+1]$ for numerical
robustness in the neighbourhood of the gimbal-lock singularity at
$\theta = \pm 90^\circ$.


% ---------------------------------------------------------------------
\subsection{Numerical Integration}
\label{sec:integration}

Two integrators are provided for the augmented state:
\begin{itemize}\setlength{\itemsep}{2pt}
  \item a hand-written fixed-step classical fourth-order Runge--Kutta
        (RK4) with quaternion re-normalisation after each step
        (zero-order hold on controller outputs over each $\Delta t$);
  \item an adaptive solver via SciPy's
        \texttt{scipy.integrate.solve\_ivp}
        (\texttt{RK45} / \texttt{DOP853} / \texttt{LSODA}), with the
        closed-loop treated as a continuous-time system and quaternion
        re-normalisation applied in post-processing.
\end{itemize}
Both integrators share the same dynamics, controller, and slosh
modules; only the time-stepping strategy differs.

%   wherever you want the section to appear.
%
% REQUIRED PACKAGES (most Overleaf templates already load these)
%       \usepackage{amsmath, amssymb}
%       \usepackage{bm}     % bold math symbols
%
% BIBLIOGRAPHY
%   Add the following entry to your .bib file so that \cite{markley2014}
%   resolves correctly:
%
%   @book{markley2014,
%     author    = {Markley, F. Landis and Crassidis, John L.},
%     title     = {Fundamentals of Spacecraft Attitude Determination and Control},
%     publisher = {Springer},
%     year      = {2014},
%     series    = {Space Technology Library},
%     volume    = {33},
%     address   = {New York, NY},
%     doi       = {10.1007/978-1-4939-0802-8}
%   }
%
%   If you are not using BibTeX, replace the \cite{markley2014} occurrences
%   below with a manual reference such as "Markley \& Crassidis (2014)".
% =====================================================================


\section{Spacecraft Attitude Dynamics and Control Model}
\label{sec:attitude-model}

This section summarises the equations of motion used in the closed-loop
attitude simulation. Sign conventions and equation numbering follow
Markley and Crassidis~\cite{markley2014} throughout.


% ---------------------------------------------------------------------
\subsection{Quaternion Convention}
\label{sec:quat-convention}

We adopt the \emph{scalar-last} unit quaternion
\begin{equation}
\bm{q} \;=\; \begin{bmatrix} \bm{q}_v \\ q_4 \end{bmatrix} \in \mathbb{R}^{4},
\qquad \bm{q}_v \in \mathbb{R}^{3},\; q_4 \in \mathbb{R},\; \|\bm{q}\| = 1,
\label{eq:scalar-last-q}
\end{equation}
matching the layout used in Eq.~(2.84) of~\cite{markley2014}. The
quaternion product is defined such that
$\bm{A}(\bm{q}\otimes\bm{p}) = \bm{A}(\bm{q})\,\bm{A}(\bm{p})$
(Eq.~2.82b of~\cite{markley2014}):
\begin{equation}
\bm{q}\otimes\bm{p} \;=\;
\begin{bmatrix}
 q_4\,\bm{p}_v + p_4\,\bm{q}_v - \bm{q}_v\times\bm{p}_v \\[2pt]
 q_4\,p_4 - \bm{q}_v^{\!\top}\bm{p}_v
\end{bmatrix},
\qquad
\bm{q}^{-1} \;=\; \begin{bmatrix} -\bm{q}_v \\ q_4 \end{bmatrix}.
\end{equation}


% ---------------------------------------------------------------------
\subsection{Quaternion Kinematics}
\label{sec:quat-kinematics}

The rate of change of the body-to-inertial quaternion under body angular
velocity $\bm{\omega}\in\mathbb{R}^{3}$ is
(Eq.~2.88 of~\cite{markley2014})
\begin{equation}
\dot{\bm{q}}
\;=\; \tfrac{1}{2}\,\bm{\Omega}(\bm{\omega})\,\bm{q}
\;=\;
\begin{bmatrix}
 \tfrac{1}{2}\bigl(q_4\,\bm{\omega} - \bm{\omega}\times\bm{q}_v\bigr) \\[3pt]
-\tfrac{1}{2}\,\bm{\omega}^{\!\top}\bm{q}_v
\end{bmatrix}.
\label{eq:quat-kin}
\end{equation}


% ---------------------------------------------------------------------
\subsection{Rigid-Body Dynamics with Reaction Wheels}
\label{sec:euler-wheels}

For a spacecraft with symmetric, positive-definite inertia tensor
$\bm{J}\in\mathbb{R}^{3\times3}$ (Eq.~3.13 of~\cite{markley2014})
and internal wheel angular momentum $\bm{h}_w^{\,b}$ expressed in the body
frame, Euler's rotational equation generalises to (Sec.~7.2
of~\cite{markley2014})
\begin{equation}
\boxed{\;
\bm{J}\,\dot{\bm{\omega}}
\;=\;
\bm{u}_{\text{body}}
\;-\;
\bm{\omega}\times\!\bigl(\bm{J}\bm{\omega} + \bm{h}_w^{\,b}\bigr).
\;}
\label{eq:euler-wheels}
\end{equation}
The total external torque on the body collects three contributions:
\begin{equation}
\bm{u}_{\text{body}} \;=\; -\bm{W}\bm{\tau}_w \;+\; \bm{\tau}_{\text{mag}}
\;+\; \bm{T}_{\!s},
\label{eq:u-body-decomp}
\end{equation}
the wheel reaction torque, the magnetic-torquer torque, and the propellant
slosh torque defined in Sections~\ref{sec:bdot} and~\ref{sec:slosh}.


% ---------------------------------------------------------------------
\subsection{Reaction Wheels and Torque Allocation}
\label{sec:wheels}

We model $n_w$ momentum wheels whose spin-axis unit vectors form the
columns of a configuration matrix
$\bm{W}\in\mathbb{R}^{3\times n_w}$. The scalar wheel momenta about
their spin axes are stacked into $\bm{h}_w\in\mathbb{R}^{n_w}$, so the
body-frame wheel momentum is
\begin{equation}
\bm{h}_w^{\,b} \;=\; \bm{W}\,\bm{h}_w.
\end{equation}
With wheel torque commands $\bm{\tau}_w\in\mathbb{R}^{n_w}$, the wheels
react $-\bm{W}\bm{\tau}_w$ onto the body and their internal momentum
evolves as
\begin{equation}
\dot{\bm{h}}_w \;=\; \bm{\tau}_w.
\end{equation}
A desired body control torque $\bm{u}_{\text{cmd}}$ from the attitude
controller is mapped to wheel commands by the minimum-norm pseudoinverse
allocation
\begin{equation}
\bm{\tau}_w \;=\; -\bm{W}^{+}\,\bm{u}_{\text{cmd}},
\qquad
\bm{W}^{+} \;=\; \bm{W}^{\!\top}\!\bigl(\bm{W}\bm{W}^{\!\top}\bigr)^{-1}.
\label{eq:wheel-alloc}
\end{equation}

\paragraph{NASA standard four-wheel pyramid.}
The simulation uses the symmetric four-wheel pyramid configuration
(Sec.~7.2 of~\cite{markley2014}). Each spin axis is tilted by
$\beta = \arccos(1/\sqrt{3}) \approx 54.7356^\circ$ from the body $+z$
axis and the four wheels are spaced $90^\circ$ apart in azimuth. This
gives equal projection onto all three body axes and provides
single-fault tolerance.


% ---------------------------------------------------------------------
\subsection{Quaternion-Feedback PD Controller}
\label{sec:pd-controller}

Given a desired attitude quaternion $\bm{q}_{\text{des}}$, the attitude
error quaternion is
\begin{equation}
\delta\bm{q} \;=\; \bm{q}_{\text{des}}^{\,-1} \otimes \bm{q}
\;=\; \bigl[\,\delta\bm{q}_v;\; \delta q_4\,\bigr]^{\!\top}.
\end{equation}
The quaternion-feedback proportional--derivative law of Sec.~7.4
of~\cite{markley2014} is used, with element-wise saturation:
\begin{equation}
\bm{u}_{\text{cmd}} \;=\;
\operatorname{sat}_{u_{\max}}\!\Bigl(
 -K_p\,\operatorname{sign}(\delta q_4)\,\delta\bm{q}_v
 \;-\; K_d\,\bm{\omega}
\Bigr),
\label{eq:pd}
\end{equation}
where $K_p, K_d > 0$ are scalar gains and
$\operatorname{sat}_{u_{\max}}(\cdot)$ clips each component to
$[-u_{\max},\, +u_{\max}]$. The factor
$\operatorname{sign}(\delta q_4)$ selects the shortest rotation and
prevents the well-known \emph{unwinding} phenomenon.


% ---------------------------------------------------------------------
\subsection{Magnetic Momentum Unloading}
\label{sec:bdot}

Reaction wheels accumulate angular momentum from disturbance torques
until they saturate. Magnetic torquers transfer that momentum from the
wheels to the geomagnetic field (Sec.~7.5 of~\cite{markley2014}). A
magnetic dipole moment $\bm{m}\in\mathbb{R}^{3}$ commanded in a field
$\bm{B}\in\mathbb{R}^{3}$ produces a body torque
\begin{equation}
\bm{\tau}_{\text{mag}} \;=\; \bm{m}\times\bm{B}.
\end{equation}
We use the cross-product (``$\bm{B}$-dot-style'') unloading law
\begin{equation}
\bm{m} \;=\;
\operatorname{sat}_{m_{\max}}\!\biggl(
\frac{K_{\text{dump}}}{\|\bm{B}\|^{2}}\,
 \bm{h}_w^{\,b} \times \bm{B}
\biggr),
\label{eq:bdot}
\end{equation}
which (before saturation) yields the body torque
\begin{equation}
\bm{\tau}_{\text{mag}}
\;=\;
-K_{\text{dump}}\,
\Bigl(\bm{h}_w^{\,b} - (\bm{h}_w^{\,b}\!\cdot\!\hat{\bm{B}})\,\hat{\bm{B}}\Bigr),
\end{equation}
i.e.\ a damping torque applied to the component of $\bm{h}_w^{\,b}$
\emph{perpendicular} to $\bm{B}$. Orbital variation of $\bm{B}$ rotates
the unloadable subspace and sweeps it through all three body axes over
each orbit. The commanded dipole is saturated element-wise to
$\bm{m}\in[-m_{\max},\,+m_{\max}]^{3}$.

For the geomagnetic field, the simulation uses a simple inclined rotating
dipole of magnitude $\sim\!30\,\mu\mathrm{T}$ as a low-Earth-orbit proxy;
in flight-grade work this is replaced by an IGRF model rotated into the
body frame.


% ---------------------------------------------------------------------
\subsection{Propellant Slosh Torque}
\label{sec:slosh}

Liquid propellant slosh is modelled as a forced, damped second-order
oscillator acting on a 3-D torque state $\bm{T}_{\!s}\in\mathbb{R}^{3}$,
driven by body rate $\bm{\omega}$ and body angular acceleration
$\dot{\bm{\omega}}$:
\begin{equation}
\boxed{\;
\ddot{\bm{T}}_{\!s}
\;=\;
 -A_s(t)\,\bm{\omega}
 \;-\; B_s(t)\,\dot{\bm{\omega}}
 \;-\; C_s(t)\,\dot{\bm{T}}_{\!s}
 \;-\; K_s(t)\,\bm{T}_{\!s}.
\;}
\label{eq:slosh}
\end{equation}
The coefficients $A_s, B_s, C_s, K_s$ are scalars (identical on all three
axes) and admit an optional sinusoidal time variation about a positive
DC offset:
\begin{equation}
f(t) \;=\; f_{\text{nom}} + f_{\text{amp}}\sin(\omega_f\,t),
\qquad f \in \{A_s,\,B_s,\,C_s,\,K_s\}.
\label{eq:slosh-coeff}
\end{equation}
Requiring $f_{\text{nom}} > |f_{\text{amp}}|$ keeps $K_s,C_s$ strictly
positive (a negative stiffness or damping would be unphysical).

The slosh torque enters Euler's equation~\eqref{eq:euler-wheels} as an
external disturbance through $\bm{T}_{\!s}$ in~\eqref{eq:u-body-decomp}.
Because $\bm{T}_{\!s}$ is a \emph{state} (not an instantaneous algebraic
function of $\dot{\bm{\omega}}$), there is no implicit loop:
$\dot{\bm{\omega}}$ is computed from~\eqref{eq:euler-wheels} using the
current $\bm{T}_{\!s}$, and the just-computed $\dot{\bm{\omega}}$ is then
substituted into~\eqref{eq:slosh}.


% ---------------------------------------------------------------------
\subsection{State Vector and Equations of Motion}
\label{sec:state}

Collecting the previous results, the full simulation state is
\begin{equation}
\bm{x}(t) \;=\;
\bigl[\,\bm{q};\; \bm{\omega};\; \bm{h}_w;\; \bm{T}_{\!s};\;
       \dot{\bm{T}}_{\!s}\,\bigr]
\;\in\; \mathbb{R}^{\,13+n_w},
\end{equation}
governed by
\begin{align}
\dot{\bm{q}} &= \tfrac{1}{2}\,\bm{\Omega}(\bm{\omega})\,\bm{q}
   && \text{(kinematics, Eq.~\ref{eq:quat-kin})} \\[2pt]
\bm{J}\,\dot{\bm{\omega}} &= \bm{u}_{\text{body}}
 \;-\; \bm{\omega}\times\!\bigl(\bm{J}\bm{\omega}+\bm{W}\bm{h}_w\bigr)
   && \text{(Euler, Eq.~\ref{eq:euler-wheels})} \\[2pt]
\dot{\bm{h}}_w &= \bm{\tau}_w \;=\; -\bm{W}^{+}\bm{u}_{\text{cmd}}
   && \text{(wheel momentum)} \\[2pt]
\ddot{\bm{T}}_{\!s} &=
   -A_s\bm{\omega} - B_s\dot{\bm{\omega}}
   - C_s\dot{\bm{T}}_{\!s} - K_s\bm{T}_{\!s}
   && \text{(slosh, Eq.~\ref{eq:slosh})}
\end{align}
with
\(
\bm{u}_{\text{body}} \;=\;
-\bm{W}\bm{\tau}_w \;+\; \bm{m}\times\bm{B} \;+\; \bm{T}_{\!s},
\)
$\bm{u}_{\text{cmd}}$ given by~\eqref{eq:pd}, and $\bm{m}$
given by~\eqref{eq:bdot}.


% ---------------------------------------------------------------------
\subsection{Euler-Angle Conversion}
\label{sec:euler-conv}

For initial-condition specification and post-processing visualisation,
3-2-1 ($Z$-$Y$-$X$) yaw--pitch--roll Euler angles $(\phi,\theta,\psi)$
are converted to and from the scalar-last quaternion using
App.~B of~\cite{markley2014}:
\begin{equation}
\bm{q} \;=\; \bm{q}_z(\psi)\otimes\bm{q}_y(\theta)\otimes\bm{q}_x(\phi).
\end{equation}
The inverse extraction, used for the angle/rate output plots, is
\begin{align}
\phi   &= \operatorname{atan2}\!\bigl(2(q_4 q_1 + q_2 q_3),\;
                                      1 - 2(q_1^{2} + q_2^{2})\bigr), \\[2pt]
\theta &= \arcsin\!\bigl(\operatorname{clip}\bigl(2(q_4 q_2 - q_1 q_3),\,
                                                  -1,\,+1\bigr)\bigr), \\[2pt]
\psi   &= \operatorname{atan2}\!\bigl(2(q_4 q_3 + q_1 q_2),\;
                                      1 - 2(q_2^{2} + q_3^{2})\bigr),
\end{align}
where the $\arcsin$ argument is clipped to $[-1,+1]$ for numerical
robustness in the neighbourhood of the gimbal-lock singularity at
$\theta = \pm 90^\circ$.


% ---------------------------------------------------------------------
\subsection{Numerical Integration}
\label{sec:integration}

Two integrators are provided for the augmented state:
\begin{itemize}\setlength{\itemsep}{2pt}
  \item a hand-written fixed-step classical fourth-order Runge--Kutta
        (RK4) with quaternion re-normalisation after each step
        (zero-order hold on controller outputs over each $\Delta t$);
  \item an adaptive solver via SciPy's
        \texttt{scipy.integrate.solve\_ivp}
        (\texttt{RK45} / \texttt{DOP853} / \texttt{LSODA}), with the
        closed-loop treated as a continuous-time system and quaternion
        re-normalisation applied in post-processing.
\end{itemize}
Both integrators share the same dynamics, controller, and slosh
modules; only the time-stepping strategy differs.
